% SUJET 1 – ÉPREUVE ANTICIPÉE (SANS CORRECTION)
\documentclass[a4paper, 12pt]{article}

\usepackage[french]{babel}
\usepackage[T1]{fontenc}
\usepackage[utf8]{inputenc}
\usepackage{amsmath, amssymb, amsfonts}
\usepackage{geometry}
\usepackage{graphicx}
\usepackage{array}
\usepackage{enumitem}
\usepackage{tikz}
\usepackage{pgfplots}
\pgfplotsset{compat=1.18}
\usepackage{tkz-tab}

\geometry{top=2cm, bottom=2cm, left=2cm, right=2cm}

% CONFIGURATION
\newcommand{\sujet}{5}
\newcommand{\classe}{Première Spécialité}
\newcommand{\duree}{2 heures}
\newcommand{\datedevoir}{Avril 2026}
\newcommand{\theme}{Épreuve anticipée – Automatismes \& Analyse}

% Cadre réponse (hauteur 2,5 cm)
\newcommand{\reponse}{\par\vspace{0.5cm}
\framebox[\linewidth]{\begin{minipage}{\dimexpr\linewidth-2\fboxsep}
\vspace{2.5cm}\hspace{0.5cm}
\end{minipage}}\vspace{0.5cm}
}

\begin{document}

\begin{center}\Large\textbf{Épreuve anticipée de mathématiques – Sujet \sujet}\end{center}

\begin{tabular}{|p{8cm}|p{8cm}|}
\hline
\textbf{Nom et Prénom :} & \textbf{Date : \datedevoir} \\
\framebox[7.5cm]{\hspace{0.1cm}\begin{minipage}{7.3cm}\vspace{1.5cm}\end{minipage}} & \\
\hline
\end{tabular}

\begin{center}\textbf{\classe – \theme}\end{center}

\textbf{Durée : \duree – Calculatrice interdite}

\section*{Partie 1 : Automatismes (6 points)}
Pour cette première partie, aucune justification n’est demandée et une seule réponse est possible par question. Reportez le numéro de la question et votre réponse dans le cadre prévu.

\begin{enumerate}[label=\textbf{Question \arabic*}, leftmargin=*]

\item On peut affirmer que :
\begin{itemize}[label=$\square$]
\item[A.] $\dfrac{5}{6} < \dfrac{6}{7}$
\item[B.] $\dfrac{5}{6} > \dfrac{6}{7}$
\item[C.] $\dfrac{5}{6} = \dfrac{6}{7}$
\item[D.] On ne peut pas comparer ces deux nombres.
\end{itemize}
\reponse

\item $1\ \text{m}^2$ correspond à :
\begin{itemize}[label=$\square$]
\item[A.] $0,0001\ \text{cm}^2$
\item[B.] $0,01\ \text{cm}^2$
\item[C.] $100\ \text{cm}^2$
\item[D.] $10\,000\ \text{cm}^2$
\end{itemize}
\reponse

\item Dans une classe de 30 élèves, 40\% sont des filles. Le nombre de filles est :
\begin{itemize}[label=$\square$]
\item[A.] 12
\item[B.] 15
\item[C.] 18
\item[D.] 20
\end{itemize}
\reponse

\item Le prix d’un article subit une augmentation de 20\,\% puis une augmentation de 10\,\%. Le taux d’évolution global est :
\begin{itemize}[label=$\square$]
\item[A.] 30\%
\item[B.] 32\%
\item[C.] 22\%
\item[D.] 35\%
\end{itemize}
\reponse

\item On donne ci‑contre la courbe représentative d’une fonction $f$.
\begin{center}
\begin{tikzpicture}[scale=0.7]
\begin{axis}[axis lines=middle, xmin=-1, xmax=5, ymin=-2, ymax=4, xtick={0,1,2,3,4,5}, ytick={-2,-1,0,1,2,3,4}, xlabel=$x$, ylabel=$y$, grid=both, grid style={gray!30}]
\addplot[domain=-0.5:4.5, samples=100, thick] {-0.5*(x-2)^2+3};
\end{axis}
\end{tikzpicture}
\end{center}
On lit graphiquement que $f(2)$ est égal à :
\begin{itemize}[label=$\square$]
\item[A.] 0
\item[B.] 1
\item[C.] 2
\item[D.] 3
\end{itemize}
\reponse

\item On lance deux fois de suite un dé truqué. La probabilité d’obtenir un 6 à chaque lancer est $0,2$. La probabilité d’obtenir  deux  6 est :
\begin{itemize}[label=$\square$]
\item[A.] $0,512$
\item[B.] $0,488$
\item[C.] $0,04$
\item[D.] $0,8$
\end{itemize}
\reponse

\item Le diagramme en barres ci‑dessous donne la répartition des élèves d’un lycée selon leur âge.
\begin{center}
\begin{tikzpicture}
\begin{axis}[ybar, ymin=0, ymax=60, symbolic x coords={16,17,18,19}, xtick=data, ytick={0,5,10,15,20,25,30,35,40,45,50,55,60}, ylabel=Effectif, xlabel=Âge, grid=major, grid style={gray!30}]
\addplot coordinates {(16,25) (17,40) (18,30) (19,5)};
\end{axis}
\end{tikzpicture}
\end{center}
La médiane de cette série est :
\begin{itemize}[label=$\square$]
\item[A.] 17
\item[B.] 17,5
\item[C.] 18
\item[D.] 18,5
\end{itemize}
\reponse

\item L’équation $x^2 = 9$ a pour solutions :
\begin{itemize}[label=$\square$]
\item[A.] $x = 3$ seulement
\item[B.] $x = -3$ seulement
\item[C.] $x = 3$ ou $x = -3$
\item[D.] $x = \sqrt{3}$ ou $x = -\sqrt{3}$
\end{itemize}
\reponse

\item Soit $f$ la fonction définie sur $\mathbb{R}$ par $f(x)=2x-4$. Le tableau de signes de $f$ est :
\begin{itemize}[label=$\square$]
\item[A.] \begin{tikzpicture}[scale=0.8]
\tkzTabInit[lgt=2,espcl=2]{$x$/1, $f(x)$/1}{$-\infty$, $2$, $+\infty$}
\tkzTabLine{ , -, 0, +, }
\end{tikzpicture}
\item[B.] \begin{tikzpicture}[scale=0.8]
\tkzTabInit[lgt=2,espcl=2]{$x$/1, $f(x)$/1}{$-\infty$, $2$, $+\infty$}
\tkzTabLine{ , +, 0, -, }
\end{tikzpicture}
\item[C.] \begin{tikzpicture}[scale=0.8]
\tkzTabInit[lgt=2,espcl=2]{$x$/1, $f(x)$/1}{$-\infty$, $0$, $+\infty$}
\tkzTabLine{ , -, 0, +, }
\end{tikzpicture}
\item[D.] \begin{tikzpicture}[scale=0.8]
\tkzTabInit[lgt=2,espcl=2]{$x$/1, $f(x)$/1}{$-\infty$, $0$, $+\infty$}
\tkzTabLine{ , +, 0, -, }
\end{tikzpicture}
\end{itemize}
\reponse

\item La forme factorisée de $(x+3)^2 - 4(x+3)$ est :
\begin{itemize}[label=$\square$]
\item[A.] $(x+3)(x-1)$
\item[B.] $(x+3)(x+1)$
\item[C.] $(x-3)(x+1)$
\item[D.] $(x-3)(x-1)$
\end{itemize}
\reponse

\item Le prix d’un article a baissé de 20\,\%. Pour retrouver son prix initial, il faut l’augmenter de :
\begin{itemize}[label=$\square$]
\item[A.] 20\%
\item[B.] 25\%
\item[C.] 30\%
\item[D.] 40\%
\end{itemize}
\reponse

\item La droite d’équation $y=3x$ passe par le point :
\begin{itemize}[label=$\square$]
\item[A.] $A(2;6)$
\item[B.] $B(3;6)$
\item[C.] $C(1;4)$
\item[D.] $D(0;3)$
\end{itemize}
\reponse

\end{enumerate}

\newpage

\section*{Partie 2 : Exercices (14 points)}

\subsection*{Exercice 1 (7 points)}
Soit $(u_n)$ la suite définie par $u_0 = 10$ et pour tout entier naturel $n$,
\[
u_{n+1} = \frac{1}{2}u_n + 2.
\]

\begin{enumerate}
    \item Calculer $u_1$ et $u_2$.
    \reponse
    \item La suite $(u_n)$ est‑elle arithmétique ? géométrique ? Justifier.
    \reponse
    \item On pose $v_n = u_n - 4$ pour tout $n\in\mathbb{N}$.
        \begin{enumerate}
            \item Montrer que $(v_n)$ est une suite géométrique dont on précisera la raison et le premier terme.
            \reponse
            \item Exprimer $v_n$ en fonction de $n$, puis $u_n$ en fonction de $n$.
            \reponse
        \end{enumerate}
    \item Étudier le sens de variation de $(u_n)$.
    \reponse
  
\end{enumerate}

\subsection*{Exercice 2 (7 points)}
On lance une balle en l’air. Sa hauteur $h(t)$ (en mètres) à l’instant $t$ (en secondes) est donnée par
\[
h(t) = -5t^2 + 20t + 1.
\]

\begin{enumerate}
    \item Quelle est la hauteur initiale de la balle ?
    \reponse
 \item Dresser le tableau de variations de $h$ sur $[0; +\infty[$.
    \reponse
    \item Déterminer l’instant $t_m$ auquel la balle atteint sa hauteur maximale.
    \reponse
    \item Calculer cette hauteur maximale.
    \reponse
   
   
\end{enumerate}

\end{document}